\section*{C5: Programare Logica}
\subsection{Teorie}
\subsubsection{Reminder: Logica de Ordinul I}

Logica de ordinul I extinde logica propozițională pentru a putea exprima fapte despre obiecte, relații și cuantificatori.

\begin{itemize}
	\item \textbf{Alfabetul:} Variabile logice $V = \{x_0, x_1, \dots\}$, constante și simboluri logice ($\perp, \rightarrow, \forall, =$).
	\item \textbf{Signatura $\sigma = (F, R, r)$:} Definește limbajul. Conține simboluri de funcții ($F$), simboluri de relații ($R$) și o funcție de aritate $r : F \cup R \to \mathbb{N}$. Constantele sunt funcții cu aritate 0.
	\item \textbf{Termenii ($T_\sigma$):} Reprezintă „obiectele”. Sunt formați recursiv din variabile și aplicarea simbolurilor de funcții (ex: $x$, $f(x, y)$, $g(c)$).
	\item \textbf{Formulele ($F_\sigma$):} Reprezintă „afirmațiile”.
	      \begin{itemize}
		      \item \textit{Formule atomice:} ecuații $t = u$ sau relații $R(t_1, \dots, t_n)$.
		      \item \textit{Formule compuse:} obținute folosind conectori și cuantificatori (ex: $\varphi \rightarrow \psi$, $\forall x \varphi$).
	      \end{itemize}
	\item \textbf{Structuri și Semantică:} O $\sigma$-structură $\mathcal{A}$ asociază fiecărui simbol o semnificație concretă (univers de discurs, funcții și relații matematice). Spunem că $\mathcal{A} \models \varphi$ dacă structura satisface formula $\varphi$.
\end{itemize}

\subsubsection{Unificarea și Teorema de Unificare}

Unificarea este procesul de găsire a unei substituții care face ca doi termeni (sau două formule) să devină identici. Este mecanismul central în programarea logică (Prolog) pentru potrivirea interogărilor cu regulile.

\subsubsection{Substituții și Unificatori}
\begin{itemize}
	\item \textbf{Substituția:} O funcție $\theta : V \to T_\sigma$ (înlocuiește variabile cu termeni). Se extinde unic la $\tilde{\theta} : T_\sigma \to T_\sigma$ prin aplicare structurală pe termeni.
	\item \textbf{Unificator:} Dată o mulțime de ecuații $\mathcal{E}$, o substituție $\theta$ este unificator dacă $\tilde{\theta}(s) = \tilde{\theta}(t)$ pentru orice ecuație $(s=t) \in \mathcal{E}$.
	\item \textbf{Cel mai general unificator (cgu / mgu):} Unificatorul $\theta$ este \textit{cgu} dacă pentru orice alt unificator $\theta'$ există o substituție $\omega$ astfel încât $\theta' = \tilde{\omega} \circ \theta$.
\end{itemize}

\subsubsection{Algoritmul de Unificare}
Primește la intrare o mulțime de ecuații $\mathcal{E}$ și returnează un \textit{cgu} sau „eșec”. Se aplică iterativ următoarele reguli până când niciuna nu mai poate fi aplicată:
\begin{enumerate}
	\item \textbf{Descompunere:} $f(s_1, \dots, s_n) = f(t_1, \dots, t_n) \implies s_1=t_1, \dots, s_n=t_n$.
	\item \textbf{Conflict (Clash):} $f(\dots) = g(\dots)$ cu $f \neq g \implies$ \textbf{eșec}.
	\item \textbf{Ștergere:} $x = x \implies$ se elimină.
	\item \textbf{Orientare:} $t = x$ (unde $t \notin V$) $\implies x = t$.
	\item \textbf{Occur-check:} $x = t$ cu $x \in Var(t)$ și $x \neq t \implies$ \textbf{eșec} (un termen nu se poate unifica cu un termen care îl conține strict).
	\item \textbf{Eliminare/Substituție:} $x = t$ cu $x \notin Var(t)$ și $x$ apare în alte ecuații $\implies$ se substituie $x$ cu $t$ peste tot în restul lui $\mathcal{E}$.
\end{enumerate}

\subsubsection{Teorema de Unificare}
\textbf{Enunț:} Algoritmul de unificare prezentat se termină mereu și produce rezultatul corect (dacă există soluție, produce un \textit{cgu} sub formă rezolvată; altfel se oprește cu „eșec”).

\subsection{Demonstrații din curs}

\subsubsection{Corectitudinea formei rezolvate}

\textbf{Propoziție:} Fie $n \in \mathbb{N}$, variabilele $z_1, \dots, z_n \in V$ (distincte) și $t_1, \dots, t_n \in T_\sigma$ astfel încât, pentru orice $i, j$, avem $z_i \notin Var(t_j)$. Atunci substituția $\theta$ definită prin $\theta(z_i) = t_i$ este \textit{cgu} pentru $\mathcal{E} = \{z_i = t_i \mid i \in \{1, \dots, n\} \}$.

\textbf{Demonstrație:}
\begin{itemize}
	\item \textit{Este unificator:} Prin definiție, aplicând $\theta$ peste $z_i$ obținem $t_i$. Cum $z_i$ nu apare în niciun $t_j$, $\tilde{\theta}(t_i) = t_i$. Deci $\tilde{\theta}(z_i) = \tilde{\theta}(t_i)$.
	\item \textit{Este cgu:} Fie $\theta'$ un alt unificator pentru $\mathcal{E}$. Definim $\omega := \theta'$. Vrem să arătăm că $\theta' = \tilde{\omega} \circ \theta$.
	      Evaluăm pe variabilele $z_i$:
	      $$(\tilde{\omega} \circ \theta)(z_i) = \tilde{\theta'}(\theta(z_i)) = \tilde{\theta'}(t_i)$$
	      Deoarece $\theta'$ unifică $z_i = t_i$, avem $\tilde{\theta'}(z_i) = \tilde{\theta'}(t_i)$, deci $\tilde{\theta'}(t_i) = \theta'(z_i)$. Astfel obținem $(\tilde{\omega} \circ \theta)(z_i) = \theta'(z_i)$.
	      Pentru orice altă variabilă $y \notin \{z_1, \dots, z_n\}$, $\theta(y) = y$, deci $(\tilde{\omega} \circ \theta)(y) = \tilde{\theta'}(y) = \theta'(y)$. Așadar, $\theta' = \tilde{\omega} \circ \theta$.
\end{itemize}

\subsubsection{Demonstrația Teoremei de Unificare}

\textbf{Demonstrație Teoremă (Terminare și Corectitudine):}

\textbf{1. Terminarea:}
Toate regulile, cu excepția ultimei (Substituția), reduc complexitatea mulțimii $\mathcal{E}$ (scad numărul de simboluri de funcție, aliniază variabilele sau elimină tautologii). Regula de substituție poate crește dimensiunea termenilor, dar ea elimină complet o variabilă din restul sistemului. Cum numărul de variabile este finit, regula 6 poate fi aplicată doar de un număr finit de ori. Așadar, algoritmul se termină.

\textbf{2. Corectitudinea pasului de Substituție (detaliere):}
Acesta este pasul critic („lăsat ca exercițiu” - afirmația că $\tilde{\theta} \circ \psi = \theta$).
Fie ecuația eliminatoare $x=t$ (cu $x \notin Var(t)$) și notăm cu $\psi$ substituția $x \mapsto t$. Restul ecuațiilor de forma $s=u$ se transformă în $\tilde{\psi}(s) = \tilde{\psi}(u)$. Trebuie să arătăm că ambele mulțimi de ecuații (înainte și după aplicarea lui $\psi$) au exact aceiași unificatori.

Fie $\theta$ un unificator pentru sistemul inițial. Deci $\theta(x) = \tilde{\theta}(t)$ și $\tilde{\theta}(s) = \tilde{\theta}(u)$.
Arătăm întâi că $\tilde{\theta} \circ \psi = \theta$:
Pentru $x$, avem $(\tilde{\theta} \circ \psi)(x) = \tilde{\theta}(t) = \theta(x)$.
Pentru $y \neq x$, avem $(\tilde{\theta} \circ \psi)(y) = \tilde{\theta}(y) = \theta(y)$.
Deci identitatea este adevărată.
Folosind identitatea de mai sus pentru ecuațiile noi $\tilde{\psi}(s) = \tilde{\psi}(u)$:
$$\tilde{\theta}(\tilde{\psi}(s)) = (\tilde{\theta} \circ \psi)(s) = \tilde{\theta}(s)$$
$$\tilde{\theta}(\tilde{\psi}(u)) = (\tilde{\theta} \circ \psi)(u) = \tilde{\theta}(u)$$
Cum $\tilde{\theta}(s) = \tilde{\theta}(u)$, rezultă că $\theta$ unifică și noua formă a ecuației.

Invers, fie $\theta$ un unificator pentru sistemul rezultat. Avem din nou $\theta(x) = \tilde{\theta}(t)$ (din ecuația menținută $x=t$). Similar se demonstrează că $\tilde{\theta} \circ \psi = \theta$.
Din faptul că $\theta$ unifică ecuația modificată $\tilde{\psi}(s) = \tilde{\psi}(u)$, avem:
$$\tilde{\theta}(\tilde{\psi}(s)) = \tilde{\theta}(\tilde{\psi}(u)) \implies (\tilde{\theta} \circ \psi)(s) = (\tilde{\theta} \circ \psi)(u) \implies \tilde{\theta}(s) = \tilde{\theta}(u)$$
Deci $\theta$ era unificator și pentru sistemul inițial, demonstrând conservarea mulțimii soluțiilor la fiecare pas.
